\documentclass[tint=true]{swisstex}

\setrunningtitle{SwissTeX \textbar{} Klassendokumentation}
\hypersetup{pdftitle={SwissTeX -- Satzsystem nach der Neuen Schweizer Schule}}

\let\code\swisscode

\begin{document}

\swisscover*{kicker={Dokumentation}, title={SwissTeX},
  subtitle={Satzsystem nach der Neuen Schweizer Schule},
  foot={Version 2.0.0 \textbar{} TeX Gyre Heros \textbar{} Raster 13,5\,pt
   \textbar{} Satzspiegel 105\,mm}}

\begin{lead}
SwissTeX leitet den gesamten Satzspiegel aus sieben Kennzahlen ab. Es gibt
keine handgesetzten Abstände. Wer eine Kennzahl ändert, ändert das System
konsistent mit. Diese Dokumentation ist mit der Klasse selbst gesetzt.
\end{lead}

\section{Invarianten}

\subsection{Die sechs Regeln}

Die Klasse ist auf sechs Aussagen gebaut. Alles andere folgt daraus.

\begin{swisstable}{Invarianten der Klasse und ihre technische Durchsetzung}
\begin{tabularx}{\textcolumn}[t]{@{}S{8mm}@{\hspace{3mm}}S{30mm}@{\hspace{3mm}}L@{}}
\toprule
\tabhead{Nr.} & \tabhead{Regel} & \tabhead{Durchsetzung} \\
\midrule
I1 & Jeder vertikale Abstand ist ein Vielfaches des Rastermasses
   & \code{\textbackslash gridskip}, \code{\textbackslash parskip} = 1 Zeile,
     Blockplatzierung mit Rasterfang \\
\addlinespace[3.5pt]
I2 & Linien und Satzblöcke stehen auf der Textachse
   & \code{\textbackslash textcolumn} als einzige Bezugsbreite; Rahmen\-elemente
     nutzen \code{\textbackslash fullmeasure} \\
\addlinespace[3.5pt]
I3 & Die Marginalspalte trägt nur Beschriftung, links Glossen, rechts Ziffern
   & \code{\textbackslash marg}, \code{\textbackslash marginlabel},
     Tabellenlegende, Zone \code{numberzone} \\
\addlinespace[3.5pt]
I4 & Eine Schriftfamilie, Unterscheidung über Schnitt, Grad, Weite
   & \code{\textbackslash setmainfont} plus \code{\textbackslash condensed} \\
\addlinespace[3.5pt]
I5 & Flattersatz mit Trennung und optischem Randausgleich
   & \code{ragged2e} und \code{microtype} \\
\addlinespace[3.5pt]
I6 & Überschriften hängen am nachfolgenden Block
   & \code{\textbackslash Needspace*}, Kopplung von Abschnitt und
     Unterabschnitt, Bindung der Tabellen \\
\bottomrule
\end{tabularx}
\end{swisstable}

\subsection{Warum das Rastermass schwierig ist}

\marg{Der schwierigste Teil der Klasse steckt in \code{\textbackslash swiss@placebox}}
Ein Block, dessen Höhe kein Vielfaches des Rastermasses ist, verschiebt allen
folgenden Satz. Das naheliegende Aufrunden der Blockhöhe genügt nicht, weil
TeX beim Anfügen einer Box Zeilenausgleich einsetzt, dessen Betrag von der
Tiefe der Vorzeile abhängt, und weil beim Seitenumbruch statt dessen
\code{\textbackslash topskip} greift.

SwissTeX packt Blöcke deshalb so um, dass sie sich wie ganze Zeilen verhalten:
Oberkante auf Struthöhe, Tiefe auf ein Vielfaches des Rastermasses aufgefüllt,
dahinter eine unsichtbare Zeile mit Standardtiefe. Damit rechnet der normale
Ausgleich mitten auf der Seite genauso wie am Seitenanfang, wo
\code{\textbackslash topskip} auf eine Rasterzeile gesetzt ist.

\section{Kennzahlen}

\subsection{Klassenoptionen}

\begin{swisstable}{Optionen und Vorgabewerte}
\begin{tabularx}{\textcolumn}[t]{@{}S{34mm}@{\hspace{3mm}}S{22mm}@{\hspace{3mm}}L@{}}
\toprule
\tabhead{Option} & \tabhead{Vorgabe} & \tabhead{Wirkung} \\
\midrule
\code{gridunit}     & 13.5pt & Basislinienraster, Durchschuss \\
\code{bodysize}     & 9.5pt  & Grundschriftgrad \\
\code{paperwidth}   & 210mm  & Papierbreite \\
\code{paperheight}  & 297mm  & Papierhöhe \\
\code{outermargin}  & 24mm   & Steg links des Rasters \\
\code{margincolumn} & 30mm   & Marginalspalte \\
\code{numberzone}   & 8mm    & rechte Zone der Marginalspalte; zwei Stellen brauchen 7\,mm \\
\code{gutter}       & 7mm    & Bund \\
\code{textcolumn}   & 105mm  & Satzspiegel, Textachse \\
\code{topmargin}    & 26mm   & oberer Steg \\
\code{gridlines}    & 51     & Rasterzeilen je Seite \\
\code{accent}       & 255,55,37 & Signalfarbe als RGB-Tripel \\
\code{band}         & leer   & Farbe der Farbfelder; leer = \code{accent} (eine Farbe, seit 2.0) \\
\code{paper}        & 234,232,208 & Grundton der Seite \\
\code{ink}          & 26,26,26 & Textfarbe \\
\code{tint}         & false  & Grundton flächig anlegen \\
\code{showgrid}     & false  & Raster sichtbar unterlegen \\
\code{rules}        & true   & Haarlinie vor Abschnitten \\
\code{lang}         & german & Sprache für Trennung \\
\code{sans}         & texgyreheros & Dateibasis der Grundschrift \\
\code{sanscondensed} & texgyreheroscn & Dateibasis der Schmalschrift \\
\code{sansname}     & TeX Gyre Heros & Familienname für den Formelsatz \\
\code{sansnameitalic} & TeX Gyre Heros Italic & kursiver Schnitt im Formelsatz \\
\code{mathfont}     & TeX Gyre DejaVu Math & Mathematikschrift \\
\code{annotationleading} & leer & Durchschuss Kolumnentitel/Kennzeichnung; leer = 2/3 Rastermass \\
\code{glossleading}      & leer & Durchschuss Glossen und Legenden; leer = 3/4 Rastermass \\
\code{footnoteleading}   & leer & Durchschuss Fussnoten und Tabellenkörper; leer = 3/4 Rastermass \\
\code{identity}     & leer   & Name einer Identitätsdatei \code{swissidentity-Name.sty} (Abschnitt~\ref{sec:identitaet}) \\
\bottomrule
\end{tabularx}
\end{swisstable}

Der rechte Steg wird nicht gesetzt, sondern berechnet: Papierbreite minus
Steg, Marginalspalte, Bund und Satzspiegel. Ebenso ergibt sich die Satzhöhe
als Rasterzeilen mal Rastermass. Damit bleibt jede Änderung konsistent.

Auch der Durchschuss der Nebentexte ist seit Version 2.0 rastergebunden statt
handgesetzt: Kolumnentitel und Kennzeichnungen laufen auf
\code{\textbackslash annotationleading} (zwei Drittel des Rastermasses),
Randglossen sowie die Legenden von Tabellen und Abbildungen auf
\code{\textbackslash glossleading}, Fussnoten und Tabellenkörper auf
\code{\textbackslash footnoteleading} (je drei Viertel) -- kommensurable
Rasterbrüche im Verhältnis 3:2 beziehungsweise 4:3. Jede der drei Optionen
lässt sich einzeln überschreiben; bleibt sie leer, leitet die Klasse den Wert
aus \code{\textbackslash gridunit} ab.

Mit Version 2.0 laufen Randglossen sowie die Legenden von Tabellen und
Abbildungen in der Grundschrift statt in der Schmalschrift: Auch bei
Müller-Brockmann selbst stand die Schmalschrift nicht für Bildlegenden,
sondern blieb dem laufenden Tabellensatz vorbehalten. Die Klasse folgt dem
jetzt; Tabellenkörper, Fussnoten und das Kolophon bleiben in der
Schmalschrift.

\marg{Eine Regel, eine Signalfarbe: \code{band} folgt \code{accent}, bis eine Identität sie ausdrücklich trennt}
Auch die Klassenoption \code{band} verliert mit Version 2.0 ihren eigenen
Vorgabewert: Er ist jetzt leer, und ein leerer Wert bedeutet
\code{band = accent} -- Bandfelder tragen dieselbe Signalfarbe wie
Auszeichnungen, statt eines zweiten, kaum unterscheidbaren Rots. Das ist die
Regel Müller-Brockmanns gegen unmotivierte Zweitfarben (im Folgenden kurz
M-B-Regel): eine Signalfarbe pro Dokument, nicht zwei. Wer ein
eigenständiges Bandrot braucht, setzt \code{band}
weiterhin -- als Klassenoption wie vor Version 2.0, oder über
\code{\textbackslash swisssetcolors\{band=...\}} in einer Identitätsdatei
oder der Präambel (Abschnitt~\ref{sec:identitaet}); der zuletzt ausgeführte
Aufruf entscheidet.

\section{Schnittstelle}

\subsection{Befehle und Umgebungen}

\marg{Die Positionsform \code{\textbackslash swisstitle\{K\}\{T\}\{U\}} gilt ab Version 2.1 als veraltet; die Slot-Form \code{\textbackslash swisstitle*} ist ihr Ersatz}
Seit Version 2.0 steht neben der klassischen Positionsform von
\code{\textbackslash swisstitle} eine Slot-Form mit Schlüssel-Wert-Syntax,
\code{\textbackslash swisstitle*\{kicker=..., title=..., subtitle=...\}}.
Beide Formen laufen durch denselben Satzkörper und erzeugen darum
identische Titelblöcke; ein leerer oder fehlender Slot lässt seinen eigenen
Abstand entfallen, statt einen leeren Block zu reservieren -- die
Positionsform kennt dasselbe Verhalten für ein leeres Argument
(\code{\{\}}). Die Positionsform bleibt in Version 2.0 unverändert nutzbar,
gilt aber ab Version 2.1 als veraltet; neue Dokumente sollten die Slot-Form
verwenden.

\begin{swisstable}{Öffentliche Schnittstelle}
\begin{tabularx}{\textcolumn}[t]{@{}S{40mm}@{\hspace{3mm}}L@{}}
\toprule
\tabhead{Aufruf} & \tabhead{Zweck} \\
\midrule
\code{\textbackslash swisstitle\{K\}\{T\}\{U\}} & Titelblock aus Kennzeichnung, Titel, Untertitel (Positionsform, ab 2.1 veraltet) \\
\addlinespace[3.5pt]
\code{\textbackslash swisstitle*\{kicker=..., title=..., subtitle=...\}} & derselbe Titelblock als Slot-Form; leere Slots ohne reservierten Abstand \\
\addlinespace[3.5pt]
\code{\textbackslash swisslogo[lines=n, axis=text|full]\{Datei\}} & Logo als Satzblock, feste Höhe in Rasterzeilen, linke Kante auf Textachse oder vollem Mass \\
\addlinespace[3.5pt]
\code{\textbackslash setrunningtitle\{...\}} & Text des Kolumnentitels \\
\addlinespace[3.5pt]
\code{\textbackslash swisssetstrings\{...\}\{...\}} & Sprache definieren
  oder ergänzen (Abschnitt~\ref{sec:sprachen}) \\
\addlinespace[3.5pt]
\code{\textbackslash swisscode\{...\}} & Schmalschrift 8,5\,pt auf
  Rastermass, für Code und Bezeichner; \code{\textbackslash code} ist ein
  Alias dafür \\
\addlinespace[3.5pt]
\code{lead} & Vorspann, ein Grad grösser als der Grundtext \\
\addlinespace[3.5pt]
\code{\textbackslash marg\{...\}} & Randglosse in der Marginalspalte \\
\addlinespace[3.5pt]
\code{swisstable\{Legende\}} & Tabellenblock, Nummer und Legende automatisch in der Marginalspalte \\
\addlinespace[3.5pt]
\code{gridblock} & beliebiger Block mit Rasterfang, auch für abgesetzte Formeln \\
\addlinespace[3.5pt]
\code{\textbackslash snaptogrid\{...\}} & Befehlsform desselben \\
\addlinespace[3.5pt]
\code{\textbackslash gridskip\{n\}} & n Rasterzeilen Raum \\
\addlinespace[3.5pt]
\code{\textbackslash tracked\{...\}} & gesperrte Versalien für Kennzeichnungen \\
\addlinespace[3.5pt]
\code{\textbackslash colophon\{...\}} & Kolophon am Fuss der letzten Seite \\
\addlinespace[3.5pt]
\code{swissfigure\{Legende\}} & Abbildungsblock, kein Gleitobjekt \\
\addlinespace[3.5pt]
\code{\textbackslash swissgraphic\{Datei\}} & Bild auf Satzspiegelbreite \\
\addlinespace[3.5pt]
\code{\textbackslash sidenote\{...\}} & nummerierte Randnotiz statt Fussnote \\
\addlinespace[3.5pt]
\code{\textbackslash gridsnap} & Rasterfang nach fremden Umgebungen \\
\addlinespace[3.5pt]
\code{\textbackslash textcolumn}, \code{\textbackslash fullmeasure}, \code{\textbackslash gridunit} & Längen für eigene Konstruktionen \\
\bottomrule
\end{tabularx}
\end{swisstable}

\subsection{Sprachen}
\label{sec:sprachen}

\marg{Unbekannter Schlüssel ist ein Fehler, unbekannte Sprache nur eine Warnung mit Rückfall auf Englisch}
Alle sichtbaren Zeichenketten der Klasse -- die Präfixe von Tabellen- und
Abbildungslegende, das Wort für ``Seite'' sowie die vier Klassifizierungsstufen
-- laufen über eine interne, nach Sprache gegliederte Tabelle statt über fest
verdrahtete Wörter. Ausgeliefert sind Deutsch und Englisch für die sieben
Schlüssel \code{table}, \code{figure}, \code{page},
\code{classification-public}, \code{classification-internal},
\code{classification-confidential} und \code{classification-strict}; die
Klassenoption \code{lang} wählt zwischen ihnen und steuert zugleich die
Trennung über \code{polyglossia}.

Eine weitere Sprache lässt sich mit \code{\textbackslash swisssetstrings}
ergänzen -- in der Präambel oder in einer eigenen Identitätsdatei, die vor
\code{\textbackslash begin\{document\}} eingebunden wird. Die Wirkung ist
global, unabhängig von der Gruppierung an der Aufrufstelle:

\begin{gridblock}
\code{\textbackslash swisssetstrings\{french\}\{table=Tableau,
figure=Illustration,}\\
\code{\hspace{1.5em}page=Page, classification-public=public, ...\}}
\end{gridblock}

Fehlt für die eingestellte Sprache ein einzelner Schlüssel, fällt die Klasse
mit einer Warnung auf die englische Zeichenkette zurück; ein Schlüssel, der
auch dort nicht existiert, ist ein Fehler. Damit lässt sich eine Sprache auch
schrittweise nachrüsten, ohne dass ein vergessener Schlüssel unbemerkt bliebe.

\subsection{Identität}
\label{sec:identitaet}

\marg{Eine Identitätsdatei nutzt ausschliesslich die öffentlichen Setzer, nie internes Patchen}
Die Klassenoption \code{identity=Name} lädt beim
\code{\textbackslash documentclass}-Aufruf eine Datei
\code{swissidentity-Name.sty}, falls gesetzt -- fehlt diese Datei, bricht
der Lauf mit einer Fehlermeldung ab, die den gesuchten Dateinamen nennt.
Eine Identitätsdatei setzt Firmen-, Farb-, Schrift- und
Klassifizierungsangaben ausschliesslich über öffentliche Befehle, die
ebenso in der Präambel eines Dokuments ohne eigene Identitätsdatei
aufrufbar sind:

\begin{swisstable}{Öffentliche Identitätsschnittstelle}
\begin{tabularx}{\textcolumn}[t]{@{}S{40mm}@{\hspace{3mm}}L@{}}
\toprule
\tabhead{Aufruf} & \tabhead{Zweck} \\
\midrule
\code{\textbackslash swissidentitymeta\{...\}} & Firmenname, Impressum, Web-Adresse \\
\addlinespace[3.5pt]
\code{\textbackslash swisssetcolors\{...\}} & Farbrollen \code{accent}, \code{paper}, \code{ink}, optional \code{band} \\
\addlinespace[3.5pt]
\code{\textbackslash newcommand\{\textbackslash swissidentityfonts\}\{...\}} & Schriftanbieter, ersetzt die Vorgabeschrift vollständig (I4, siehe unten) \\
\addlinespace[3.5pt]
\code{\textbackslash swisslogofiles\{...\}} & Logo-Datei für den Umschlag \\
\addlinespace[3.5pt]
\code{\textbackslash swissclassifications\{...\}} & geordnetes Klassifizierungsvokabular \\
\addlinespace[3.5pt]
\code{\textbackslash swissfootformat\{...\}} & Vorlage der Fusszeile (Abschnitt~\ref{sec:fusszeile}) \\
\addlinespace[3.5pt]
\code{\textbackslash swisscovervariant\{Name\}\{...\}} & benannte Umschlagvariante registrieren (Abschnitt~\ref{sec:umschlag}) \\
\addlinespace[3.5pt]
\code{\textbackslash swissmeta\{...\}} & dokumentseitige Angaben:
  \code{docid}, \code{version}, \code{date}, \code{classification}, \code{client} \\
\bottomrule
\end{tabularx}
\end{swisstable}

\marg{Der \code{colophon}-Schlüssel von \code{\textbackslash swisslogofiles} wird abgelegt, aber von \code{\textbackslash colophon} in Version 2.0 noch nicht gelesen -- ein offen dokumentierter Randfall, kein stiller Fehler}
Alle Setzer wirken global, unabhängig von der Gruppierung an der
Aufrufstelle -- wie \code{\textbackslash swisssetstrings} in
Abschnitt~\ref{sec:sprachen}. \code{\textbackslash swissmeta} ist der
einzige dokumentseitige Befehl der Gruppe: Er trägt die Angaben eines
einzelnen Dokuments, \code{date} setzt sich ohne eigene Angabe auf
\code{\textbackslash today}. Die Klassifizierung ist ein Wert aus dem mit
\code{\textbackslash swissclassifications} gesetzten Vokabular (Vorgabe:
\code{public}, \code{internal}, \code{confidential}, \code{strict}, siehe
Abschnitt~\ref{sec:sprachen}); ein Wert ausserhalb dieses Vokabulars ist ein
Fehler, kein stiller Rückfall. Ohne \code{identity=} bleibt die Klasse bei
ihrer Vorgabe. Fusszeile, Umschlag und Kolophon lesen diese Angaben.

\subsubsection{Vollständigkeit ist Pflicht des Anbieters (I4)}

Ersetzt eine Identität die Schrift über
\code{\textbackslash newcommand\{\textbackslash swissidentityfonts\}\{...\}},
muss sie darin -- neben \code{\textbackslash setmainfont} -- auch
\code{\textbackslash renewfontfamily\textbackslash condensed\{...\}}
aufrufen (die Schmalschrift-Familie ist bereits deklariert, ein
\code{\textbackslash newfontfamily} scheiterte an dieser Stelle); ersetzt
sie auch die Mathematik, gilt dieselbe Pflicht für einen eigenen
\code{\textbackslash setmathfont}-Aufruf. Die Klasse prüft nur, ob
\code{\textbackslash swissidentityfonts} definiert ist, nicht, ob es
vollständig ist: Eine Identität, die nur \code{\textbackslash setmainfont}
austauscht, aber die Schmalschrift vergisst, baut anstandslos -- mit einer
stillen Mischung zweier Schriftfamilien, genau dem I4-Verstoss, den die
Klasse eigentlich ausschliesst. Durchgesetzt wird diese Vollständigkeit erst
nachträglich, am fertigen Dokument: swisscheck A16
(Abschnitt~\ref{sec:pruefung}) vergleicht jede eingebettete Schrift gegen
die im Kennzahlen-Sidecar deklarierten Familien. Die Referenzidentität
\code{acme} erfüllt die Pflicht vollständig und wird genau darauf geprüft.

\subsubsection{Beispiel: die Referenzidentität acme}

\code{acme} ist ein neutraler Platzhaltername, keine reale Firma. Ein
Auszug aus \code{swissidentity-acme.sty}:

\begin{gridblock}
\code{\textbackslash swissidentitymeta\{company=Acme AG, web=acme.example\}}\\
\code{\textbackslash swisssetcolors\{accent=200,16,46, paper=245,244,238, ink=26,26,26\}}\\
\code{\textbackslash newcommand\{\textbackslash swissidentityfonts\}\{}\\
\code{\hspace{1.5em}\textbackslash setmainfont\{SwissTeX Grotesk\}}\\
\code{\hspace{1.5em}\textbackslash renewfontfamily\textbackslash condensed\{SwissTeX Grotesk Condensed\}}\\
\code{\hspace{1.5em}\textbackslash setmathfont\{TeX Gyre DejaVu Math\}\}}\\
\code{\textbackslash swisslogofiles\{cover=acme-logo.pdf\}}\\
\code{\textbackslash swissclassifications\{public,internal,confidential,strict\}}\\
\code{\textbackslash swissfootformat\{\textbackslash meta\{docid\} · \textbackslash meta\{version\} · \textbackslash meta\{date\}\}}
\end{gridblock}

Ein Dokument, das diese Identität nutzt, trägt nur noch seine eigenen
Angaben:

\begin{gridblock}
\code{\textbackslash documentclass[identity=acme]\{swisstex\}}\\
\code{\textbackslash swissmeta\{docid=TR-2026-014, version=1.2, classification=internal\}}
\end{gridblock}

\code{acme-demo.tex}/\code{acme-demo.pdf} setzt diese Identität tatsächlich
um -- über jeden öffentlichen Setzer dieses Abschnitts hinweg -- und ist
damit das lauffähige Gegenstück zu dieser Beschreibung.

\subsection{Umschlag und Varianten}
\label{sec:umschlag}

\code{\textbackslash swisscover*\{...\}} legt die randabfallende
Umschlagseite mit einer Schlüssel-Wert-Liste an und steht seit Version 2.0
vollständig auf dem Seitenraster: weder das Farbband noch das
Überformatzeichen tragen noch eine Seitenanteilskoordinate.

\begin{swisstable}{Schlüssel von \code{\textbackslash swisscover*}}
\begin{tabularx}{\textcolumn}[t]{@{}S{26mm}@{\hspace{3mm}}L@{}}
\toprule
\tabhead{Schlüssel} & \tabhead{Wirkung} \\
\midrule
\code{kicker} & Kennzeichnung oben links; ab \code{internal} rechtsbündig mit der Klassifizierungsmarke \\
\addlinespace[3.5pt]
\code{title} & Haupttitel, \code{\textbackslash swissdisplay\{4\}} \\
\addlinespace[3.5pt]
\code{subtitle} & Untertitel in der Signalfarbe, \code{\textbackslash swissdisplay\{1\}} \\
\addlinespace[3.5pt]
\code{glyph} & Überformatzeichen; ohne diesen Schlüssel erscheint keins -- seit 2.0 kein Vorgabezeichen mehr \\
\addlinespace[3.5pt]
\code{glyphline} & Rasterzeile der Zeichen-Grundlinie (Vorgabe 30) \\
\addlinespace[3.5pt]
\code{glyphlines} & Zeichenhöhe in Rasterzeilen (Vorgabe 34) \\
\addlinespace[3.5pt]
\code{band} & Farbband als \code{start/lines} in Rasterzeilen (Vorgabe 38/9) \\
\addlinespace[3.5pt]
\code{logo} & Logodatei; Vorgabe aus \code{\textbackslash swisslogofiles} \\
\addlinespace[3.5pt]
\code{foot} & Fusstext links im Farbband \\
\addlinespace[3.5pt]
\code{client} & Auftraggeber, Teil der rechtsbündigen Metazeile \\
\addlinespace[3.5pt]
\code{docid} & Dokumentkennung; Vorgabe aus \code{\textbackslash swissmeta} \\
\addlinespace[3.5pt]
\code{date} & Datum; Vorgabe aus \code{\textbackslash swissmeta}, letztlich \code{\textbackslash today} \\
\addlinespace[3.5pt]
\code{variant} & registrierte Umschlagvariante anwenden; explizite Schlüssel überschreiben sie \\
\bottomrule
\end{tabularx}
\end{swisstable}

\marg{Kein Seitenanteil mehr: Band und Zeichen hängen an Rasterzeilen}
Das Farbband platziert seine Oberkante auf Rasterzeile \code{start} und
seine Höhe auf \code{lines} Rasterzeilen, über die volle Seitenbreite; die
Vorgabe \code{38/9} gilt, solange \code{band} nicht gesetzt wird. Beide
Teilwerte müssen nichtnegative Ganzzahlen sein; reicht die Spanne über das
Seitenraster hinaus, warnt die Klasse, bricht aber nicht ab. Das
Überformatzeichen folgt derselben Rasterlogik: \code{glyphline} hängt seine
Grundlinie an eine Rasterzeile, \code{glyphlines} bestimmt seine Höhe in
Rasterzeilen, und seine linke Kante liegt auf der Textachse -- dieselbe
Koordinatenkonvention wie das Rasterkontroll-Overlay
(\code{showgrid=true}). Anders als vor Version 2.0 gibt es dafür kein
Vorgabezeichen mehr: Ohne \code{glyph=} erscheint gar keins. Ein Zeichen ist
damit eine bewusste, identitätstragende Entscheidung, kein automatisches
Ornament.

\marg{Varianten legen Positionen fest, Dokumente füllen Inhalt}
\code{\textbackslash swisscovervariant\{Name\}\{Vorgaben\}} hinterlegt eine
Schlüssel-Wert-Liste unter einem Variantennamen -- aus einer
Identitätsdatei heraus (Abschnitt~\ref{sec:identitaet}) oder direkt in der
Präambel, mit derselben globalen Wirkung wie jeder andere Setzer dieses
Handbuchs. \code{\textbackslash swisscover*\{variant=Name, ...\}} wendet die
hinterlegten Vorgaben an; jeder im selben Aufruf ausdrücklich gesetzte
weitere Schlüssel überschreibt seine Entsprechung aus der Variante --
unabhängig davon, ob er vor oder nach \code{variant=...} in der Liste
steht. Die Referenzidentität \code{acme} registriert so die Variante
\code{report}:

\begin{gridblock}
\code{\textbackslash swisscovervariant\{report\}\{band=38/9, glyphline=30\}}
\end{gridblock}

-- hier deckungsgleich mit der Klassenvorgabe, aber unter einem eigenen
Namen buchbar, damit ein Dokument \code{variant=report} schreibt statt jeden
Wert einzeln zu wiederholen. Ein Vokabular mit tatsächlich unterschiedlichen
Varianten (etwa für ein Angebot oder einen internen Bericht) ist ebenso
möglich; die Klasse schliesst kein Vokabular, sie liefert nur den
Mechanismus.

\marg{Die alte Positionsform warnt bereits in Version 2.0, nicht erst ab 2.1 wie bei \code{\textbackslash swisstitle}}
Die alte Positionsform \code{\textbackslash
swisscover\{Kennzeichnung\}\{Titel\}\{Untertitel\}\{Fuss\}} bleibt nutzbar,
meldet aber schon jetzt eine Warnung beim Bau, nicht erst ab einer
künftigen Version. Dieses Handbuch setzt seinen eigenen Umschlag seit
dieser Überarbeitung deshalb mit der Slot-Form.

\subsection{Fusszeile}
\label{sec:fusszeile}

\marg{Kein Logo im Kolumnentitel: eine wiederholte Marke informiert einmal, jede weitere Wiederholung ist nur noch Werbung}
Kopf- und Fusszeile teilen sich dieselbe Rastergeometrie und dieselbe
Zurückhaltung. Der Kolumnentitel trägt links die Pagina im Grundschriftgrad
in der Signalfarbe und rechts den mit \code{\textbackslash setrunningtitle}
gesetzten Lauftitel in der Schmalschrift auf Kennzeichnungsgrad -- kein
Logo; das bleibt Umschlag und Kolophon vorbehalten
(Abschnitt~\ref{sec:umschlag}). \code{headsep} bindet den Kopf an
\code{2\textbackslash gridunit}, \code{footskip} den Fuss an
\code{3\textbackslash gridunit} -- beide Grundlinien liegen auf derselben
verlängerten Rasterzeile wie der Grundtext, geprüft von swisscheck A10 und
A11 (Abschnitt~\ref{sec:pruefung}).

Die Fusszeile trägt zwei getrennte Zonen: links, in der Marginalspalte, die
Klassifizierungsmarke, rechts, ab der Textachse, eine Metazeile aus
Dokumentangaben. Die Marke erscheint nur ab der Klassifizierungsstufe
\code{internal} aufwärts -- Schweigen ist der öffentliche Zustand, keine
sichtbare Markierung für die niedrigste Stufe. Die Metazeile kommt aus
\code{\textbackslash swissfootformat\{...\}} (Abschnitt~\ref{sec:identitaet}):
eine mit \code{\textbackslash meta\{key\}} formulierte Vorlage, die erst
beim tatsächlichen Fusssatz an die Dokumentangaben aus
\code{\textbackslash swissmeta} gebunden wird, sodass ein späteres
\code{\textbackslash swissmeta} in der Identitätsdatei oder im Dokument
selbst noch greift. Ohne eigenes \code{\textbackslash swissfootformat}
bleibt die Metazeile leer; ein einzelner nie gesetzter
\code{\textbackslash meta\{key\}}-Schlüssel liefert dabei ebenfalls nur
Leere statt eines Fehlers.

\subsection{Tabellen}

Innerhalb von \code{swisstable} wird die Tabelle direkt geschrieben. Das
\code{[t]} ist zwingend, sonst fluchtet die Legende auf die Tabellenmitte
statt auf die Kopfzeile:

\begin{gridblock}
\code{\textbackslash begin\{swisstable\}\{Legende\}}\\
\code{\textbackslash begin\{tabularx\}\{\textbackslash textcolumn\}[t]\{@\{\}S\{30mm\}@\{\textbackslash hspace\{4mm\}\}L@\{\}\}}\\
\code{\hspace{1.5em}...}\\
\code{\textbackslash end\{tabularx\}}\\
\code{\textbackslash end\{swisstable\}}
\end{gridblock}

\marg{\code{tabularx} lässt sich nicht in eine eigene Umgebung wickeln, weil es seinen Rumpf bis zum wörtlichen Ende einliest}
Bereitgestellt sind die Spaltentypen \code{L} (dehnbar, Flattersatz),
\code{S\{Breite\}} (fest, Flattersatz) und \code{R\{Breite\}}
(fest, rechtsbündig) sowie \code{\textbackslash tabhead} für die Kopfzeile.
Senkrechte Linien sind nicht vorgesehen.

\section{Fremde Umgebungen}

\subsection{Automatischer Rasterfang}

\marg{Der Fang muss nach dem Schlusscode der Umgebung greifen, nicht davor}
Umgebungen aus LaTeX und aus Paketen bringen eigene Abstände mit. Für
\code{quote}, \code{quotation}, \code{verbatim}, \code{description} und die
abgesetzten Formelumgebungen von \code{amsmath} setzt die Klasse deshalb
\code{\textbackslash gridsnap} automatisch dahinter. Der Fang liest
\code{\textbackslash pagetotal}, füllt auf die nächste Rasterzeile auf und
setzt die Tiefe für den folgenden Zeilenausgleich.

Damit brauchen abgesetzte Formeln kein \code{gridblock} mehr. Für eigene
oder fremde Umgebungen, die noch nicht erfasst sind, genügt ein Aufruf von
\code{\textbackslash gridsnap} dahinter oder ein Eintrag über
\code{\textbackslash gridsnapenv\{name\}}.

\subsection{Listen, Fussnoten, dritte Ebene}

Aufzählung, Nummerierung, verschachtelte Listen und Beschreibungslisten
stehen auf dem Raster; die Marken sitzen auf der Textachse und hängen nicht
in den Bund. Fussnoten laufen in der Schmalschrift, die Trennlinie steht auf
der Textachse statt auf vierzig Prozent der Spaltenbreite. Alternativ setzt
\code{\textbackslash sidenote} eine nummerierte Notiz in die Marginalspalte.
Die dritte Gliederungsebene ist unnummeriert und läuft als fetter Vorspann in
den Absatz hinein.

\section{Schaugrade und Farbe}

\subsection{Anzeigengrade}

\marg{Auch Schaugrade bleiben am Raster: der Durchschuss ist ein Vielfaches, der Grad folgt daraus}
\code{\textbackslash swissdisplay\{n\}} setzt einen Grad mit n Rasterzeilen
Durchschuss. Der Schriftgrad wird daraus abgeleitet, ab drei Zeilen läuft die
Schrift enger, wie im Schaugebrauch der Grotesk üblich.

\begin{gridblock}
{\swissdisplay{3}\bfseries Grotesk}\par
{\swissdisplay{2}Satzspiegel}\par
{\swissdisplay{1}Basislinienraster}
\end{gridblock}

Die Abschnittsziffer steht jetzt auf zwei Rasterzeilen in der
Marginalspalte. Das ist das auffälligste Gliederungsmittel der Neuen
Schweizer Schule: Ordnung über Grösse und Lage statt über zusätzliche
Auszeichnung.

\subsection{Farbfelder}

Farbfelder sind Rahmensatz. Sie laufen über das volle Mass, ihre Höhe bleibt
über dieselbe Blockplatzierung am Raster, und die Spaltenregel gilt in ihnen
nicht.

\begin{swissband}
{\swissdisplay{2}\bfseries Kunsthalle}\par\vspace{6pt}
\code{\textbackslash begin\{swissband\}} setzt ein Feld in der Bandfarbe mit
weissem Satz. \code{\textbackslash swisscover} legt eine randabfallende
Umschlagseite an, \code{tint=true} den Grundton über alle Seiten.
\end{swissband}

Die Palette ist voreingestellt auf einen warmen Grundton und ein Signalrot;
das Farbfeld trägt seit Version 2.0 standardmässig dieselbe Signalfarbe
statt eines eigenen, helleren Bandrots. Alle vier Rollen lassen sich als
Klassenoption setzen: \code{paper}, \code{accent}, \code{band}, \code{ink}.

\subsection{Blindzeile über dem Farbfeld}

\marg{Satzblock und Farbfläche sind zwei verschiedene Dinge und brauchen zwei verschiedene Abstände}
Ein Satzblock beginnt mit einer Textzeile, deren Grundlinie ins Raster
gehört; über ihr steht die Oberlänge als optischer Abstand. Eine Farbfläche
hat keine Oberlänge, ihre Kante ist die Kante. Wird sie wie ein Satzblock
gesetzt, stösst sie fast an die Unterlängen der Vorzeile.

Müller-Brockmann legt zwischen Satz und Feld eine Blindzeile: gemessen von
der Unterlänge der Vorzeile bis zur Feldkante, und von der Feldkante bis zur
Oberlänge der Folgezeile. \code{\textbackslash swissbandlead} zählt diese
Blindzeilen, voreingestellt auf eine. Ganzzahlig ist dann nicht das Feld,
sondern das Modul aus Blindzeile, Feld und Blindzeile. Der verbleibende
Unterschied im sichtbaren Weissraum entspricht der Differenz aus Ober- und
Unterlänge und lässt sich mit ganzen Rasterzeilen nicht wegrechnen.

\section{Prüfung}
\label{sec:pruefung}

\subsection{swisscheck}

Die Klasse kann Regeln nur erzwingen, solange die Schnittstelle benutzt wird.
Handgesetzte Abstände oder ein vergessenes \code{[t]} fallen erst am Ergebnis
auf. Das beiliegende Skript misst deshalb die fertige PDF-Datei:

\begin{gridblock}
\code{python3 swisscheck.py bericht.pdf -{}-tex bericht.tex}
\end{gridblock}

Geprüft werden Achsentreue der Linien, Rasterbindung des Grundtexts gegen die
Satzoberkante, rechte Satzkante, Marginalspalte und Bund, Legendenflucht der
Tabellen, Überschriften am Seitenfuss, die beiden Zonen der Marginalspalte
sowie Zeilenanfänge ohne Einzug. Die Kennzahlen lassen sich per
Schalter an ein abweichendes Raster anpassen.

\marg{Ein häufiger Stolperstein bei eigenen Prüfungen: PDF rechnet in DTP-Punkt, TeX in TeX-Punkt}
Wer eigene Messungen anstellt, muss zwischen TeX-Punkt und DTP-Punkt
umrechnen. Ohne diesen Faktor erscheint jedes korrekte Raster als langsam
wegdriftend, mit etwa 0,05\,pt je Zeile.

\subsection{Kennzahlen-Sidecar}

\marg{Die Sidecar-Datei prüft rückwirkend auch A1 bis A9: sobald sie vorliegt, gelten deklarierte Kennzahlen statt Kommandozeilen-Vorgaben}
Seit Version 2.0 schreibt die Klasse \code{\textbackslash AtEndDocument}
eine zweite, kleine Datei neben die PDF: \code{<jobname>.swisscheck}, eine
Zeile \code{schlüssel=wert} je Kennzahl -- Rastermass, die drei
Nebentext-Durchschüsse, Kopf- und Fusszonen, Steg- und Spaltenmasse,
Farbrollen, Schriftfamilien, Klassifizierung, Dokumentkennung.
\code{swisscheck.py} findet diese Sidecar-Datei automatisch neben der
geprüften PDF (oder nimmt sie über \code{-{}-params Datei} entgegen) und
prüft dann gegen die deklarierten Kennzahlen dieses Dokuments statt gegen
feste Kommandozeilen-Annahmen -- ein Dokument mit eigenem
\code{outermargin=} oder eine Identität mit eigenem \code{accent=} bleibt
ohne manuelle Schalter prüfbar. Das gilt rückwirkend auch für A1 bis A9:
Liegt eine Sidecar-Datei vor, lesen auch die ursprünglichen neun Prüfungen
ihre Kennzahlen aus ihr, nicht aus den Kommandozeilen-Vorgaben. Fehlt die
Datei -- etwa bei einer älteren, vor Version 2.0 gebauten PDF --, fällt
swisscheck für A1 bis A9 auf seine bisherigen Vorgaben zurück und
überspringt die sechs neuen Prüfungen sauber, ohne einen Verstoss zu
melden.

Sechs neue Prüfungen kommen damit hinzu: die Fusszeile auf dem
Rasterfusspunkt mit ihren beiden Zonen (A10, Abschnitt~\ref{sec:fusszeile}),
der Kolumnentitel mit der Pagina im Grundschriftgrad und einer bildfreien
Kopfzone (A11), der Umschlag mit dem Farbband auf Rasterzeilen und
Schriftgraden aus der Anzeigenskala (A12, Abschnitt~\ref{sec:umschlag}), die
Kommensurabilität des Nebentext-Durchschusses in der Marginalspalte gegen
die drei deklarierten Kennzahlen (A13), die Unterscheidbarkeit eines
eigenständigen Bandtons von der Signalfarbe (A15) sowie das
Schriftinventar der eingebetteten Schriften gegen die deklarierten
Familien (A16). A14, die Prüfung für zweisprachigen Satz, ist für eine
spätere Ausbaustufe vorgesehen und heute nicht implementiert.

\subsection{Grenzen}

Die Rasterbindung gilt für Grundtext. Tabellenkörper und Fussnoten laufen
bewusst in der Schmalschrift auf eigenem, engerem Durchschuss; Randglossen
sowie die Legenden von Tabellen und Abbildungen laufen zwar in der
Grundschrift, aber ebenfalls auf einem eigenen, vom Fliesstext
unterschiedenen Durchschuss. Alle stehen innerhalb ihres Blocks auf eigenem
Mass; nach aussen bleibt der Block rastertreu. Gleitobjekte sind
nicht vorgesehen, weil sie mit einem Basislinienraster und der Marginalspalte
schlecht zusammengehen: Abbildungen stehen dort, wo sie geschrieben sind. Der
Fussnotenbereich steht am Seitenfuss und folgt einem eigenen Mass.

\colophon{SwissTeX 2.0.0, gesetzt mit sich selbst. TeX Gyre Heros als
Univers-Analogon, 9,5/13,5\,pt im Flattersatz. Satzspiegel 105\,mm,
Marginalspalte 30\,mm, 51 Rasterzeilen je Seite. XeLaTeX.}

\end{document}
